\documentclass[12pt]{article}
\usepackage[utf8]{inputenc}
\usepackage[T1]{fontenc}
\usepackage{hyperref}
\title{Yongkom Language Resources for K-12 Curriculum Development}
\author{Compiled by Claude Sonnet 5.5 (claude-sonnet-5-5) for Andr\'as Kornai}
\date{}
\begin{document}
\maketitle

\section{Introduction}
Yongkom (also Yonggom or Yongom; ISO 639-3 \texttt{yon}, Glottocode yong1280) is an Ok language of the Western Province of Papua New Guinea, spoken along the Fly and Tedi (Alice) rivers and toward Lake Murray. It also extends into Papua province in Indonesia. Speaker counts in the sources conflict. Ethnologue gives about 6,000, a Bible preface gives about 4,000 in Papua New Guinea, and one linguistic site gives about 20,000 including Indonesia. Search results turned up very few educational resources.

\section{Primary Educational Resources}
No primers, textbooks or school programmes were found.

\section{Linguistic and Reference Materials}
\begin{itemize}
\item Overview of the language: \url{https://en.wikipedia.org/wiki/Yonggom_language}
\item Endangered Languages Project entry: \url{https://endangeredlanguages.com/node/98502}
\item Yonggom fieldwork materials (Australian research data catalogue): \url{https://researchdata.edu.au/yonggom-fieldwork-materials-new-guinea/1408847}
\item Counting-system data: \url{https://lingweb.eva.mpg.de/channumerals/Yonggom.htm}
\end{itemize}

\section{Religious and Community Resources}
A Yongkom New Testament (title page ``God ye weng karadmo'') is published on eBible: \url{https://ebible.org/yon/FRT01.htm}

\section{Digital and Audio Resources}
No audio, video or software resources were found in the searches.

\section{Curriculum Development Context}
The documented written material is limited to the New Testament and linguistic data. A curriculum in Yongkom would require new literacy materials, ideally developed with the community and with SIL-style orthography work. The sources I found do not confirm the current orthography status.
\end{document}